% TOP_PROBLEM: {"id": 3172, "title": "Hamiltonicity of Cayley graphs", "category_id": 3, "category": {"id": 3, "name": "graph_theory", "display_name": "Graph Theory", "description": "Problems involving graphs, networks, and their properties.", "slug": "graph-theory", "order_index": 3, "created_at": "2026-07-31T15:26:25.671Z"}, "status": "open", "problem_number": "OPG-2095", "set_id": 12, "statusQualification": "The terse source is interpreted in its standard finite connected undirected form, with a generating inverse-closed set and group order at least three."}
% FIRST_POSTED: 2026-10-01
% ATTEMPT_STATUS: unresolved
% UnsolvedMath ID 3172; source catalogue code OPG-2095
% !TeX program = lualatex
% GPT-6 Astra Ultra research attempt; independent partial work, 2026-10-01.
\documentclass[11pt,a4paper]{article}
\usepackage[margin=25mm]{geometry}
\usepackage{fontspec}
\setmainfont{lmroman10-regular.otf}[BoldFont=lmroman10-bold.otf,
 ItalicFont=lmroman10-italic.otf,BoldItalicFont=lmroman10-bolditalic.otf]
\usepackage{amsmath,amssymb,mathtools}
\usepackage{unicode-math}
\setmathfont{latinmodern-math.otf}
\AtBeginDocument{\let\setminus\smallsetminus}
\usepackage{xurl,hyperref}
\hypersetup{colorlinks=true,urlcolor=blue}
\setcounter{secnumdepth}{0}
\setlength{\parindent}{0pt}
\setlength{\parskip}{5pt}
\setlength{\emergencystretch}{3em}
\begin{document}
\section*{Hamiltonicity of Cayley graphs}\label{3172}
{\small UnsolvedMath ID 3172 · OPG-2095 · Graph Theory · OpenGarden}
\subsection{Definitions and mathematical statement}
Let $\Gamma$ be a finite group with identity $1$ and order
$N\geq3$. Choose an inverse-closed generating set
$S\subseteq\Gamma\setminus\{1\}$, so $S^{-1}=S$ and
$\langle S\rangle=\Gamma$. The undirected Cayley graph
$\operatorname{Cay}(\Gamma,S)$ has vertex set $\Gamma$ and
an edge between distinct $g,h$ precisely when
$g^{-1}h\in S$. Inverse closure makes it undirected,
excluding $1$ removes loops, and the generating condition
makes the graph connected. Repeated descriptions of an edge
give only one edge.

The Hamiltonicity question is whether every such graph
contains a Hamilton cycle. Explicitly, for every pair
$(\Gamma,S)$ there should exist an enumeration
\[
 g_1,g_2,\ldots,g_N
\]
of the group elements such that $g_i^{-1}g_{i+1}\in S$
for all indices modulo $N$.

The enumeration must use every group element exactly once
before returning to its start. The generating set is part
of the input and cannot be replaced by a more convenient
one. The group is not assumed abelian. This records the
standard finite connected interpretation of the source's
short question: an arbitrary disconnected Cayley graph
cannot have one spanning cycle, and orders one and two
do not contain a simple cycle under the convention here.
Directed Cayley graphs without inverse-closed sets constitute
a different problem and are not included.
\subsection{Short English statement}
Does every finite connected undirected Cayley graph with at least three vertices have a Hamilton cycle?
\subsection{Research attempt: lifting a quotient cycle with generating voltage}
\paragraph{Scope and literature check.}
GPT-6 Astra Ultra, 1 October 2026. The graph is undirected and the
specified generating set must be respected. Morris and Wilk [S3]
prove the conjecture for group orders $kp$ with $k<48$ and $p$
prime, with the order-two exception excluded here already. Morris'
primary paper [S4] treats another family and uses quotient lifting.
The construction below proves the relevant lifting lemma directly and
checks explicit cyclic and dihedral cases. It does not infer a
solution for arbitrary groups from those substantial restricted results.

\paragraph{1. A cyclic normal-subgroup lifting certificate.}
Let $A$ be a cyclic normal subgroup of $\Gamma$. Suppose a word
$s_1\cdots s_m$, with every $s_i\in S$, traces a Hamilton cycle in
the quotient Cayley graph on $\Gamma/A$, starting at the identity
coset, where $m=|\Gamma/A|\ge3$. Put
\[
 p_0=1,\quad p_i=s_1\cdots s_i,\quad w=p_m\in A.
\]
The quotient condition says $Ap_0,\ldots,Ap_{m-1}$ are all distinct
and exhaust the cosets. Assume in addition that $w$ generates $A$.
Repeat the word $|A|$ times in the original Cayley graph. During
block $t$, its vertices before successive steps are
\[
 w^t p_i,\qquad 0\le t<|A|,\quad 0\le i<m.
\]
Each step is right multiplication by its prescribed $s_i$, so it
is an allowed edge under the canonical convention. Different $i$
give different cosets. For fixed $i$, equality
$w^t p_i=w^{t'}p_i$ implies $w^t=w^{t'}$ and hence $t=t'$ in
the chosen range. Thus the list contains every group element exactly
once. Its final product is $w^{|A|}=1$, so the walk closes and
is a Hamilton cycle. No centrality of $A$ is used: left translating
a path by $w^t$ preserves all of its right-multiplication edge labels.

\paragraph{2. The exact obstruction when voltage fails.}
If $w$ has order $d<|A|$, the same repeated word closes after only
$d$ blocks and covers exactly $dm$ vertices. Starting in the other
left translates by elements of $A$ produces $|A|/d$ disjoint lifted
cycles. Indeed their vertices in each coset are the disjoint sets
$a\langle w\rangle p_i$. Thus a quotient Hamilton cycle by itself
produces a cycle cover, not necessarily one Hamilton cycle upstairs.
The missing condition is that its voltage $w$ generate the whole
normal subgroup, or that one have an additional operation joining the
lifted components. Normality and connectedness alone do not make this
particular voltage generate $A$.

\paragraph{3. Two direct families using the actual generators.}
If some $s\in S$ has order $|\Gamma|=N$, then
$1,s,s^2,\ldots,s^{N-1},1$ is immediately a Hamilton cycle. This
covers cyclic groups only when the given set contains a generator
of the full cyclic group; such a generator is not asserted to occur
in every generating set.

For the dihedral group
\[
 D_{2m}=\langle r,s\mid r^m=s^2=1,\ srs=r^{-1}\rangle,
 \qquad m\ge3,
\]
assume $\{r,r^{-1},s\}\subseteq S$. The sequence
\[
 1,r,\ldots,r^{m-1},r^{m-1}s,r^{m-2}s,\ldots,s,1
\]
is a Hamilton cycle. The first layer advances by right multiplication
by $r$, the layer transitions use $s$, and the second layer also
advances by $r$ because $(r^j s)r=r^{j-1}s$. The two cosets of
$\langle r\rangle$ are disjoint, so all $2m$ vertices occur once.
Additional generators add edges and do not invalidate the certificate.
This proves the claim for this specified family of generating sets,
not for every generating set of every dihedral group.

\paragraph{Verification and first remaining gap.}
The script [S5] implements dihedral multiplication explicitly and checks
the cycle for $3\le m\le20$. It also verifies quotient lifting
in $\mathbb Z/12\mathbb Z$ with normal subgroup $3\mathbb Z/12\mathbb Z$
and the length-three quotient word consisting of three unit steps.
Those checks confirm the conventions; the lifting proof is uniform.
For a general group one may lack a suitable cyclic normal subgroup,
and even when one exists, a generating-voltage quotient cycle is not
constructed here. These are actual unresolved hypotheses of the method.

\subsection{Final status}
UNRESOLVED. The quotient-lifting criterion and the stated explicit
families are proved. The general connected Cayley-graph assertion is
not established. Completion estimate: no defensible global percentage.

\subsection{Sources}
\begin{itemize}
\item[S1] ULAM.AI and the UnsolvedMath Contributors, supplied problems.json, record 3172 (OPG-2095), OpenGarden collection. Catalogue identity and source transcription; CC BY 4.0.
\url{https://huggingface.co/datasets/ulamai/UnsolvedMath}.
\item[S2] Open Problem Garden, Hamiltonicity of Cayley graphs. Original problem and discussion; the mathematical formulation below is expanded from the supplied source record.
\url{http://www.openproblemgarden.org/op/hamiltonicity_of_cayley_graphs}.
\item[S3] D. Witte Morris and K. Wilk, \emph{Cayley graphs of order $kp$ are hamiltonian for $k<48$}, arXiv:1805.00149 (2018). Checked 1 October 2026.
\url{https://arxiv.org/abs/1805.00149}.
\item[S4] D. Witte Morris, \emph{On hamiltonian cycles in Cayley graphs of order $pqrs$}, arXiv:2107.14787 (2021). Checked 1 October 2026.
\url{https://arxiv.org/abs/2107.14787}.
\item[S5] GPT-6 Astra Ultra, exact group-word checks, 1 October 2026: \texttt{scripts/check\_attempt\_batch02\_connectivity\_2026\_10\_01.py} in this repository; run with Python 3.
\end{itemize}
\end{document}
